% Reference-aware execution of the window targets, part of section 3.6.
\paragraph{Admit and reclaim without evicting live references.}
The Master offers hot blocks from the first $L$ layers in slack order,
using stable request IDs to break ties. Existing copies can be shared;
new copies require both budget and physical space, reserved before
transfer. Under pressure, only zero-reference objects with no in-flight
reader or writer and valid DRAM backing enter the LRU eviction queue.
If free space and safe victims are insufficient, the Peer increment is
denied and the request retains its DRAM path. Normal inference does not
wait for a Peer grant, and priority does not revoke a live reference.

A successful grant holds a request reference until that LLM request
ends, including a response that hands control to an external tool.
Completion decrements shared references; zero references make a copy
eligible for reclamation, not immediate deletion. Growing or changed
hot sets require fresh admission. Previously granted objects remain
protected, so the current candidate count is a target rather than a
guarantee of complete Peer coverage or a bound on retained old copies.

When the budget shrinks, new allocations stop above the target and
safe zero-reference copies are reclaimed first. Referenced copies drain
as requests finish; reducing $L$ likewise does not revoke their pins.
Physical usage can therefore exceed a newly reduced target, and memory
is returned to foreground inference only after release completes.
Increasing $L$ queues additional candidates subject to the same checks.
Copying, committed, and reclaiming bytes remain charged until physical
release. Window targets never imply instant copies or instant free space.
